\chapter*{怎样读这一卷}
\addcontentsline{toc}{chapter}{怎样读这一卷}

本卷从二八〇年孙吴降服写至五八九年隋军入建康、陈亡。西晋、东晋、十六国、南朝、北朝以及东西魏、北齐、北周和隋都在这里进入年序；二八〇与五八九只是最高政权的转换点，并不表示户籍、寺院、军镇或地方社会在同一天完成重组。

\begin{center}
\begin{tikzpicture}[x=.036cm,y=1cm,font=\small]
  \fill[paper] (280,-.72) rectangle (589,.9);
  \draw[-{Stealth[length=3mm]},line width=.8pt,color=source] (280,0) -- (593,0);

  \foreach \year in {280,317,383,420,494,534,589} {
    \draw[timeline!85,line width=.55pt] (\year,-.12) -- (\year,.12);
  }
  \node[align=center,font=\scriptsize] at (280,.42) {280\\西晋统一};
  \node[align=center,font=\scriptsize] at (317,.8) {317\\东晋立};
  \node[align=center,font=\scriptsize] at (383,.42) {383\\淝水};
  \node[align=center,font=\scriptsize] at (420,.8) {420\\刘宋立};
  \node[align=center,font=\scriptsize] at (494,.42) {494\\迁都洛阳};
  \node[align=center,font=\scriptsize] at (534,.8) {534\\东西魏分立};
  \node[align=center,font=\scriptsize] at (589,.42) {589\\隋灭陈};

  \draw[dynasty,line width=2.4pt] (280,-.36) -- (316,-.36);
  \node[below,font=\scriptsize] at (298,-.39) {西晋};
  \draw[timeline,line width=2.4pt] (317,-.36) -- (420,-.36);
  \node[below,font=\scriptsize] at (368,-.39) {东晋与十六国};
  \draw[dynasty,line width=2.4pt] (420,-.36) -- (534,-.36);
  \node[below,font=\scriptsize] at (477,-.39) {南朝与北魏};
  \draw[timeline,line width=2.4pt] (534,-.36) -- (589,-.36);
  \node[below,font=\scriptsize] at (561,-.39) {东西魏至隋陈};
\end{tikzpicture}

\smallskip
\small\textit{图\quad 两晋与南北朝的阅读时间轴，280--589年：按《晋书》、南北朝正史、《隋书》与《资治通鉴》的年序标出本册各阶段的进入点。它只帮助定位并行政权和章节转换，不以单一朝代线替代十六国、地方军镇、迁徙、制度执行或宗教网络的不同节奏。}
\end{center}

\phantomsection\label{fig:jin-northern-southern-timeline}


先循年号、都城、迁徙、继承和战争把并行的政权放回各自时段，再用专题追问道路、军镇、侨置、土地、宗教与行政怎样改变选择。不要把“分裂—统一”当成自动因果：请区分长期的资源与社会条件、眼前触发、让行动得以发生的组织条件，以及结果怎样在别处制造新的压力。

《晋书》、南北朝正史与《资治通鉴》提供年序骨架；墓志、造像记、佛道文献、河西和吐鲁番文书、城址与道路遗存显示不同地方的痕迹。十六国的年号和疆域、南渡规模、族群称谓、制度执行和战争数字都须按材料的保存与立场保留边界；一方墓志或一处遗址也不能代表全部社会。

信息框帮助暂停而不打断叙事：\texttt{event} 指向首次深描，\texttt{turningpoint} 标出改变后续路径的节点，\texttt{causalchain} 分开条件、触发和影响；\texttt{textwindow} 限定一段原文所能照亮的范围；\texttt{sourcenote}、\texttt{debate} 说明证据与分歧。读完框后回到前后事件，检验它究竟改变了哪一条时间线。
